\begin{abstract}
We formalize the Birth--death Moran process on a finite graph and the ``if'' direction of the
isothermal theorem of Lieberman, Hauert and Nowak (2005): on a connected regular graph, a
mutant population of size $k$ with fitness $r\neq 1$ takes over the whole population with
probability $(1-r^{-k})/(1-r^{-n})$, and with probability $k/n$ when $r=1$. The complete graph
gives Moran's formula (1958). As in the voter project, eventual probabilities are suprema of
increasing finite-time probabilities; no path-space measure is used.
\end{abstract}

\section{The model}
\begin{definition}[Birth--death Moran kernel]\label{def:moran}
\lean{Moran.Config,Moran.mutants,Moran.fitness,Moran.totalFitness,Moran.target,Moran.pairDist,Moran.moranKernel}\leanok
Mutants have fitness $r>0$ and residents fitness $1$. A step chooses a parent $u$ with
probability proportional to its fitness and an offspring position $w$ uniformly among the
neighbours of $u$ (or $w=u$ if $u$ is isolated); the individual at $w$ is replaced by a copy
of the parent.
\end{definition}
\begin{lemma}[The kernel is well defined]\label{lem:well-defined}
\lean{Moran.target_nonneg,Moran.target_sum_one,Moran.totalFitness_pos,Moran.pair_weight_nonneg,Moran.pair_weight_sum_one}\leanok\uses{def:moran}
The weights of (parent, offspring position) pairs are nonnegative and sum to one.
\end{lemma}
\begin{proof}\leanok
Sum over the offspring position first: each parent's placement weights sum to one.
The remaining sum of fitnesses over the total fitness is one.
\end{proof}

\section{Invariance on regular graphs}
\begin{lemma}[Birth mass is $r$ times death mass]\label{lem:birth-death}
\lean{Moran.orientedCut,Moran.orientedCut_symm,Moran.birthMass,Moran.deathMass,Moran.typed_mass,Moran.birth_eq_r_death}\leanok\uses{lem:well-defined}
On a $d$-regular graph, the probability that a step adds a mutant equals $r$ times the
probability that it removes one: both are proportional to the number of edges between the two
types, counted from either side.
\end{lemma}
\begin{proof}\leanok
For $d>0$ both masses are the oriented cut divided by $Fd$, weighted by the parent's fitness;
the cut is symmetric under swapping endpoints. For $d=0$ both masses vanish.
\end{proof}
\begin{theorem}[Invariant potentials]\label{thm:invariance}
\lean{Moran.apply_affine,Moran.moran_potential_invariant,Moran.moran_mutants_invariant}\leanok\uses{lem:birth-death}
On a regular graph, $(1/r)^{\#\text{mutants}}$ is preserved in expectation by one step; when
$r=1$, so is the number of mutants.
\end{theorem}
\begin{proof}\leanok
A step changes the number of mutants by $\pm1$ or not at all, so its expectation is affine in
the birth and death masses, and these cancel by the previous lemma.
\end{proof}

\section{Absorption on connected graphs}
\begin{theorem}[Fixation or extinction]\label{thm:absorption}
\lean{Moran.unfixed,Moran.moran_constant,Moran.unfixed_step,Moran.unfixed_access,Moran.moran_unfixed_tendsto}\leanok\uses{lem:well-defined}
On a connected graph, the probability that neither type has taken over tends to zero.
\end{theorem}
\begin{proof}\leanok
Constant configurations are absorbing. From a mixed configuration, a mutant adjacent to a
resident (which exists by connectivity) reproduces onto it with positive probability; by
induction on the number of residents, fixation is reachable. Apply the shared finite-chain
absorption theorem.
\end{proof}

\section{Fixation probabilities}
\begin{lemma}[Fixation as a limit]\label{lem:fixation-limit}
\lean{Moran.allMutant,Moran.fixation,Moran.fixation_mono,Moran.fixation_tendsto,Moran.fixation_eq_of_invariant}\leanok\uses{thm:absorption}
The finite-time fixation probabilities increase to the fixation probability. If an invariant
observable agrees with the fixation indicator on constant configurations and lies between it
and the indicator plus the unfixed indicator elsewhere, it equals the fixation probability.
\end{lemma}
\begin{proof}\leanok
The all-mutant configuration is absorbing, so the finite-time probabilities are monotone and
bounded. The observable differs from the fixation indicator by at most the unfixed indicator,
whose expectation tends to zero.
\end{proof}
\begin{theorem}[Isothermal theorem]\label{thm:isothermal}
\lean{Moran.moranPsi,Moran.moranPsi_invariant,Moran.moranPsi_sandwich,Moran.isothermal}\leanok\uses{thm:invariance,lem:fixation-limit}
On a connected regular graph with $r\neq1$, the fixation probability from $k$ mutants is
$(1-r^{-k})/(1-r^{-n})$.
\end{theorem}
\begin{proof}\leanok
Apply the previous lemma to $\psi=(1-(1/r)^{\#\text{mutants}})/(1-(1/r)^n)$, invariant by the
invariance theorem, equal to $1$ and $0$ on the two constant configurations and in $[0,1]$
otherwise.
\end{proof}
\begin{theorem}[Neutral case]\label{thm:neutral}
\lean{Moran.neutralPsi,Moran.neutralPsi_invariant,Moran.neutralPsi_sandwich,Moran.isothermal_neutral}\leanok\uses{thm:invariance,lem:fixation-limit}
On a connected regular graph with $r=1$, the fixation probability is $k/n$.
\end{theorem}
\begin{proof}\leanok
The same argument with $\psi=\#\text{mutants}/n$.
\end{proof}
\begin{theorem}[Moran's formula]\label{thm:moran}
\lean{Moran.moran_formula}\leanok\uses{thm:isothermal}
On the complete graph, the fixation probability from $k$ mutants is $(1-r^{-k})/(1-r^{-n})$.
\end{theorem}
\begin{proof}\leanok
The complete graph is connected and $(n-1)$-regular.
\end{proof}
The converse of the isothermal theorem is false (Galanis, Göbel, Goldberg, Lapinskas and
Richerby, Proposition~12), so only this direction is formalized.
